\chapter{Finite-Temperature Corrections to the Effective Potential}

\section{Corrections to Scalar Field Theory}

In the zero-temperature case, the one-loop contribution to the effective
potential came from the sum
\begin{equation}
	i\sum_{n=1}^{\infty}\int\frac{d^4k}{(2\pi)^4}\,\frac{1}{2n}
	\left[\frac{\frac{1}{2}\lambda\bar\varphi^2}{k^2+i\varepsilon}\right]^n
	= -\frac{i}{2}\int\frac{d^4k}{(2\pi)^4}
	\ln\left[1 - \frac{\frac{1}{2}\lambda\bar\varphi^2}{k^2+i\varepsilon}\right].
\end{equation}
At finite temperature, the replacement of the continuous energy integral
by a discrete Matsubara sum gives
\begin{equation}
	V_{\text{eff}}^{(1)\beta}(\bar\varphi)
	= U(\bar\varphi) - \frac{i}{2}\cdot\frac{1}{-i\beta}
	\sum_{n=-\infty}^{\infty}\int\frac{d^3k}{(2\pi)^3}
	\ln\left[1 - \frac{\lambda\bar\varphi^2}
		{2(\omega_n^2-k^2)}\right],
\end{equation}
\begin{equation}
	= U(\bar\varphi) + \frac{1}{2\beta}
	\sum_{n=-\infty}^{\infty}\int\frac{d^3k}{(2\pi)^3}
	\ln\left[1 - \frac{\lambda\bar\varphi^2}
		{2(\omega_n^2-k^2)}\right].
\end{equation}
For Bose fields, the Matsubara frequencies satisfy
$-i\beta\omega_n = 2\pi n$, giving $\omega_n = 2\pi in/\beta$ and
$\omega_n^2 = -4\pi^2n^2/\beta^2$. Substituting,
\begin{equation}
	V_{\text{eff}}^{(1)\beta}(\bar\varphi)
	= U(\bar\varphi) + \frac{1}{2\beta}
	\sum_{n=-\infty}^{\infty}\int\frac{d^3k}{(2\pi)^3}
	\ln\left[1 + \frac{\lambda\bar\varphi^2}
		{2\left(\frac{4\pi^2}{\beta^2}n^2+k^2\right)}\right].
\end{equation}
Writing the argument of the logarithm as a ratio,
\begin{multline}
	V_{\text{eff}}^{(1)\beta}(\bar\varphi)
	= U(\bar\varphi) + \frac{1}{2\beta}
	\sum_{n=-\infty}^{\infty}\int\frac{d^3k}{(2\pi)^3} \\
	\times\left(\ln\left[\frac{4\pi^2}{\beta^2}n^2
		+ k^2 + \frac{\lambda\bar\varphi^2}{2}\right]
	- \ln\left[\frac{4\pi^2}{\beta^2}n^2+k^2\right]\right),
\end{multline}
where the second term contributes only a field-independent constant and
is dropped. Denoting $E_k^2 = k^2 + \frac{\lambda\bar\varphi^2}{2}$,
\begin{equation}
	V_{\text{eff}}^{(1)\beta}(\bar\varphi)
	= U(\bar\varphi) + \frac{1}{2\beta}
	\int\frac{d^3k}{(2\pi)^3}
	\sum_{n=-\infty}^{\infty}
	\ln\left[\frac{4\pi^2}{\beta^2}n^2+E_k^2\right]
	+ \text{const.}
\end{equation}

\subsection*{Evaluation of the Matsubara Sum}

We evaluate the sum
\begin{equation}
	X = \sum_{n=-\infty}^{\infty}
	\ln\left[\frac{4\pi^2}{\beta^2}n^2+E_k^2\right]
\end{equation}
by differentiating with respect to $E_k$. Setting $a^2 = 4\pi^2/\beta^2$
and $y = \beta E_k/2\pi$,
\begin{align}
	\frac{dX}{dE_k}
	 & = \sum_{n=-\infty}^{\infty}\frac{2E_k}{a^2n^2+E_k^2} \notag        \\
	 & = \frac{2}{E_k} + \frac{4}{a}\sum_{n=1}^{\infty}\frac{y}{y^2+n^2},
\end{align}
where we separated the $n=0$ term. Using the identity
\begin{equation}
	\sum_{n=1}^{\infty}\frac{y}{y^2+n^2}
	= -\frac{1}{2y} + \frac{\pi}{2}\coth\pi y
	= -\frac{1}{2y} + \frac{\pi}{2}
	+ \frac{\pi e^{-2\pi y}}{1-e^{-2\pi y}},
\end{equation}
\begin{align}
	\frac{dX}{dE_k}
	 & = \frac{2}{E_k} + \frac{2\beta}{\pi}
	\left(-\frac{1}{2y}+\frac{\pi}{2}
	+\frac{\pi e^{-2\pi y}}{1-e^{-2\pi y}}\right)
	\notag                                                       \\
	 & = \frac{2}{E_k} - \frac{\beta}{\pi y}
	+ \beta + \frac{2\beta e^{-2\pi y}}{1-e^{-2\pi y}}
	\notag                                                       \\
	 & = \beta + \frac{2\beta e^{-\beta E_k}}{1-e^{-\beta E_k}},
\end{align}
where in the last step we used $y = \beta E_k/2\pi$ so that
$\beta/(\pi y) = 2/E_k$, cancelling the first term. Integrating back,
\begin{equation}
	X = \int\left(\beta + \frac{2\beta e^{-\beta E_k}}
	{1-e^{-\beta E_k}}\right)dE_k
	= \beta E_k + 2\ln\left(1-e^{-\beta E_k}\right).
\end{equation}
Substituting back,
\begin{align}
	V_{\text{eff}}^{(1)\beta}(\bar\varphi)
	 & = U(\bar\varphi) + \frac{1}{2\beta}
	\int\frac{d^3k}{(2\pi)^3}
	\left(\beta E_k + 2\ln\left(1-e^{-\beta E_k}\right)\right)
	\notag                                          \\
	 & = U(\bar\varphi) + \int\frac{d^3k}{(2\pi)^3}
	\left[\frac{E_k}{2}
		+ \frac{1}{\beta}\ln\left(1-e^{-\beta E_k}\right)\right].
\end{align}
The first term inside the integral is the zero-point energy. Using the
identity
\begin{equation}
	i\int_{-\infty}^{\infty}\frac{dx}{2\pi}
	\ln(-x^2+y^2-i\varepsilon) = y + \text{const.},
\end{equation}
and rotating to Euclidean space via $k^0\to ik^0_E$, $-k_0^2+E_k^2
	= k_E^2 + \frac{\lambda\bar\varphi^2}{2}$, the zero-temperature part
is recovered and we can write
\begin{equation}
	V_{\text{eff}}^{(1)\beta}(\bar\varphi)
	= \underbrace{U(\bar\varphi) + \frac{1}{2}
		\int\frac{d^4k_E}{(2\pi)^4}
		\ln\!\left[1+\frac{\lambda\bar\varphi^2}{2k_E^2}\right]}_{
		\text{zero-temperature result}}
	+ \underbrace{\frac{1}{2\pi^2\beta}
		\int_0^{\infty}dk\,k^2
		\ln\!\left(1-e^{-\beta E_k}\right)}_{
		\text{finite-temperature contribution}}.
\end{equation}
Substituting $x = k\beta$, $dk = dx/\beta$, $E_k =
		(k^2+\lambda\bar\varphi^2/2)^{1/2}$,
\begin{equation}
	\Delta V_{\text{eff}}^{(1)}(\bar\varphi)
	= \frac{1}{2\pi^2\beta^4}\int_0^{\infty}dx\,x^2
	\ln\!\left[1-\exp\!\left[{-\left(x^2
				+\frac{\lambda\bar\varphi^2\beta^2}{2}\right)^{1/2}}\right]\right].
\end{equation}
Defining $y^2 = \lambda\bar\varphi^2\beta^2/2$, the finite-temperature
correction can be written compactly as
\begin{equation}
	V_{\text{eff}}^{(1)}(\bar\varphi)
	= V_{\text{eff}}^{(1)\,T=0}(\bar\varphi)
	+ \frac{1}{2\pi^2\beta^4}\,I(y),
\end{equation}
where
\begin{equation}
	I(y) = \int_0^{\infty}dx\,x^2
	\ln\!\left[1-e^{-(x^2+y^2)^{1/2}}\right].
\end{equation}
The zero-temperature renormalization counterterms also render this
finite-temperature contribution finite, so no additional renormalization
is needed.

\subsection*{High-Temperature Expansion}

In the high-temperature limit $\bar\varphi\beta\ll 1$, we expand $I(y)$
in powers of $y^2$. Setting $f(y) = \ln(1-e^{-(x^2+y^2)^{1/2}})$,
\begin{align}
	f(y) & = f(0) + y^2 f'(0) + \mathcal{O}(y^4) \notag                       \\
	     & = \ln(1-e^{-x}) + y^2\left[\frac{e^{-(x^2+y^2)^{1/2}}}
		                            {1-e^{-(x^2+y^2)^{1/2}}}
		                            \cdot\frac{1}{2(x^2+y^2)^{1/2}}\right]_{y=0}
	+ \mathcal{O}(y^4) \notag                                                 \\
	     & = \ln(1-e^{-x})
	+ \frac{y^2}{2x}\cdot\frac{e^{-x}}{1-e^{-x}}
	+ \mathcal{O}(y^4).
\end{align}
Substituting into the integral,
\begin{equation}
	I(y) = \int_0^{\infty}dx\,x^2\ln(1-e^{-x})
	+ \frac{y^2}{2}\int_0^{\infty}dx\,
	\frac{xe^{-x}}{1-e^{-x}} + \mathcal{O}(y^4).
\end{equation}
For the first term, using $\ln(1-e^{-x})
	= -\sum_{n=1}^{\infty}e^{-nx}/n$,
\begin{equation}
	\int_0^{\infty}dx\,x^2\ln(1-e^{-x})
	= -\sum_{n=1}^{\infty}\frac{1}{n}
	\int_0^{\infty}dx\,x^2 e^{-nx}.
\end{equation}
Substituting $y = nx$,
\begin{equation}
	\int_0^{\infty}dx\,x^2 e^{-nx}
	= \frac{1}{n^3}\int_0^{\infty}dy\,y^2 e^{-y}
	= \frac{\Gamma(3)}{n^3} = \frac{2}{n^3},
\end{equation}
so the first term gives
\begin{equation}
	-\sum_{n=1}^{\infty}\frac{1}{n}\cdot\frac{2}{n^3}
	= -2\sum_{n=1}^{\infty}\frac{1}{n^4}
	= -2\zeta(4) = -2\cdot\frac{\pi^4}{90}
	= -\frac{\pi^4}{45}.
\end{equation}
For the second term, using $e^{-x}/(1-e^{-x})
	= \sum_{n=0}^{\infty}e^{-(n+1)x}$,
\begin{equation}
	\int_0^{\infty}dx\,\frac{xe^{-x}}{1-e^{-x}}
	= \sum_{n=1}^{\infty}\int_0^{\infty}dx\,x\,e^{-nx}.
\end{equation}
Substituting $y = nx$,
\begin{equation}
	\sum_{n=1}^{\infty}\frac{1}{n^2}
	\int_0^{\infty}dy\,y\,e^{-y}
	= \sum_{n=1}^{\infty}\frac{\Gamma(2)}{n^2}
	= \zeta(2) = \frac{\pi^2}{6}.
\end{equation}
Combining,
\begin{equation}
	\boxed{I(y) = -\frac{\pi^4}{45}
		+ \frac{\pi^2 y^2}{12} + \mathcal{O}(y^4).}
\end{equation}
Substituting back with $y^2 = \lambda\bar\varphi^2\beta^2/2$,
\begin{align}
	V_{\text{eff}}^{(1)\beta}(\bar\varphi)
	 & = V_{\text{eff}}^{(1)\,T=0}(\bar\varphi)
	+ \frac{1}{2\pi^2\beta^4}
	\left(-\frac{\pi^4}{45}
	+ \frac{\pi^2}{12}\cdot\frac{\lambda\bar\varphi^2\beta^2}{2}
	+ \mathcal{O}(\bar\varphi^4)\right) \notag        \\
	 & = \boxed{V_{\text{eff}}^{(1)\,T=0}(\bar\varphi)
		     - \frac{\pi^2}{90\beta^4}
		     + \frac{\lambda\bar\varphi^2}{24\beta^2}
		     + \mathcal{O}(\bar\varphi^4).}
\end{align}
The first correction $-\pi^2/(90\beta^4) = -\pi^2 T^4/90$ is the
familiar Stefan-Boltzmann term, the free energy density of a massless
scalar gas. The second correction $\lambda\bar\varphi^2/(24\beta^2)
	= \lambda T^2\bar\varphi^2/24$ is a temperature-dependent mass term,
positive and proportional to $T^2$. This term is what drives symmetry
restoration at high temperature: it adds a positive contribution to the
effective mass of the field, pushing the minimum of the effective
potential back toward $\bar\varphi = 0$ as $T$ increases.

\section{Corrections to the GSW Model}

The finite-temperature correction to the effective potential in the GSW
model receives contributions from the two $W$ bosons and the $Z$ boson,
each with three polarizations. Using the field-dependent masses
$m(W_{1,2}) = \frac{1}{2}g\bar\varphi$ and
$m(Z) = \frac{1}{2}(g^2+g'^2)^{1/2}\bar\varphi$, the correction is
\begin{equation}
	\Delta V_{\text{eff}}^{(1)\beta}(\bar\varphi)
	= \frac{3}{2\pi^2\beta^4}\left\{
	2I\!\left(\tfrac{1}{2}g\bar\varphi\beta\right)
	+ I\!\left[\tfrac{1}{2}(g^2+g'^2)^{1/2}\bar\varphi\beta\right]
	\right\}.
\end{equation}
Expanding each term using $I(y) = -\pi^4/45 + \pi^2 y^2/12 +
	\mathcal{O}(y^4)$, the $W$ boson contribution gives
\begin{align}
	2I\!\left(\tfrac{1}{2}g\bar\varphi\beta\right)
	 & = 2\left[-\frac{\pi^4}{45}
		      + \frac{\pi^2}{12}\cdot\frac{1}{4}g^2\bar\varphi^2\beta^2
		      + \mathcal{O}(\bar\varphi^4)\right] \notag \\
	 & = -\frac{2\pi^4}{45}
	+ \frac{\pi^2}{24}g^2\beta^2\bar\varphi^2
	+ \mathcal{O}(\bar\varphi^4),
\end{align}
and the $Z$ boson contribution gives
\begin{align}
	I\!\left[\tfrac{1}{2}(g^2+g'^2)^{1/2}\bar\varphi\beta\right]
	 & = -\frac{\pi^4}{45}
	+ \frac{\pi^2}{12}\cdot\frac{1}{4}(g^2+g'^2)
	\bar\varphi^2\beta^2
	+ \mathcal{O}(\bar\varphi^4) \notag \\
	 & = -\frac{\pi^4}{45}
	+ \frac{\pi^2}{48}(g^2+g'^2)\bar\varphi^2\beta^2
	+ \mathcal{O}(\bar\varphi^4).
\end{align}
Summing these,
\begin{align}
	\Delta V_{\text{eff}}^{(1)\beta}(\bar\varphi)
	 & = \frac{3}{2\pi^2\beta^4}\left[
		                            -\frac{2\pi^4}{45}
		                            + \frac{\pi^2}{24}g^2\beta^2\bar\varphi^2
		                            - \frac{\pi^4}{45}
		                            + \frac{\pi^2}{48}(g^2+g'^2)\bar\varphi^2\beta^2
		                            + \mathcal{O}(\bar\varphi^4)\right] \notag  \\
	 & = \frac{3}{2\pi^2\beta^4}\left[
		                            -\frac{3\pi^4}{45}
		                            + \frac{\pi^2\beta^2\bar\varphi^2}{24}
		                            \left(g^2 + \frac{g^2}{2} + \frac{g'^2}{2}\right)
		                            + \mathcal{O}(\bar\varphi^4)\right] \notag \\
	 & = \frac{3}{2\pi^2\beta^4}\left[
		                            -\frac{\pi^4}{15}
		                            + \frac{\pi^2}{48}\beta^2(3g^2+g'^2)\bar\varphi^2
		                            + \mathcal{O}(\bar\varphi^4)\right] \notag \\
	 & = -\frac{\pi^2}{10\beta^4}
	+ \frac{1}{32\beta^2}(3g^2+g'^2)\bar\varphi^2
	+ \mathcal{O}(\bar\varphi^4).
\end{align}
The full finite-temperature effective potential in the GSW model is
therefore
\begin{equation}
	\boxed{V_{\text{eff}}^{(1)\beta}(\bar\varphi)
		= V_{\text{eff}}^{(1)\,T=0}(\bar\varphi)
		- \frac{\pi^2}{10\beta^4}
		+ \frac{1}{32\beta^2}(3g^2+g'^2)\bar\varphi^2
		+ \mathcal{O}(\bar\varphi^4).}
\end{equation}
As in the scalar case, the $-\pi^2/(10\beta^4)$ term is the
Stefan-Boltzmann free energy of the gauge bosons, and the
$(3g^2+g'^2)\bar\varphi^2/(32\beta^2)$ term is a positive
temperature-dependent mass correction that drives symmetry restoration
at high temperature.

\section{Corrections to the Minimal SU(5) Model}

For the minimal SU(5) GUT, the gauge bosons acquire a field-dependent
mass $m = \sqrt{25/8}\,g\bar\varphi$ from the Higgs mechanism. The
finite-temperature correction to the effective potential is
\begin{align}
	V_{\text{eff}}^{(1)\beta}(\bar\varphi)
	 & = V_{\text{eff}}^{(1)\,T=0}(\bar\varphi)
	+ \frac{18}{\pi^2\beta^4}
	\,I\!\left(\sqrt{\tfrac{25}{8}}\,g\bar\varphi\beta\right)
	\notag                                                               \\
	 & = V_{\text{eff}}^{(1)\,T=0}(\bar\varphi)
	+ \frac{18}{\pi^2\beta^4}\left[
		                         -\frac{\pi^4}{45}
		                         + \frac{\pi^2}{12}\cdot\frac{25}{8}
		                         \,g^2\bar\varphi^2\beta^2
		                         + \mathcal{O}(\bar\varphi^4)\right] \notag \\
	 & = V_{\text{eff}}^{(1)\,T=0}(\bar\varphi)
	- \frac{2\pi^2}{5\beta^4}
	+ \frac{18}{\pi^2\beta^4}\cdot\frac{25}{8}
	\,g^2\bar\varphi^2\beta^2\cdot\frac{\pi^2}{12}
	+ \mathcal{O}(\bar\varphi^4),
\end{align}
where the factor of 18 counts the degrees of freedom of the heavy gauge
bosons in the SU(5) model. Collecting terms,
\begin{equation}
	\boxed{V_{\text{eff}}^{(1)\beta}(\bar\varphi)
		= V_{\text{eff}}^{(1)\,T=0}(\bar\varphi)
		- \frac{2\pi^2}{5\beta^4}
		+ \frac{75}{16\beta^2}\,g^2\bar\varphi^2
		+ \mathcal{O}(\bar\varphi^4).}
\end{equation}
The structure is the same as in the scalar and GSW cases: a
Stefan-Boltzmann term $-2\pi^2/(5\beta^4)$ from the free energy of the
gauge bosons, and a positive temperature-dependent mass term
$75g^2\bar\varphi^2/(16\beta^2)$ that drives symmetry restoration. The
larger coefficient relative to the GSW case reflects the larger gauge
group and the correspondingly greater number of degrees of freedom.





